\documentclass[10pt]{article}\usepackage[]{graphicx}\usepackage[]{xcolor}
% maxwidth is the original width if it is less than linewidth
% otherwise use linewidth (to make sure the graphics do not exceed the margin)
\makeatletter
\def\maxwidth{ %
  \ifdim\Gin@nat@width>\linewidth
    \linewidth
  \else
    \Gin@nat@width
  \fi
}
\makeatother

\definecolor{fgcolor}{rgb}{0.345, 0.345, 0.345}
\newcommand{\hlnum}[1]{\textcolor[rgb]{0.686,0.059,0.569}{#1}}%
\newcommand{\hlsng}[1]{\textcolor[rgb]{0.192,0.494,0.8}{#1}}%
\newcommand{\hlcom}[1]{\textcolor[rgb]{0.678,0.584,0.686}{\textit{#1}}}%
\newcommand{\hlopt}[1]{\textcolor[rgb]{0,0,0}{#1}}%
\newcommand{\hldef}[1]{\textcolor[rgb]{0.345,0.345,0.345}{#1}}%
\newcommand{\hlkwa}[1]{\textcolor[rgb]{0.161,0.373,0.58}{\textbf{#1}}}%
\newcommand{\hlkwb}[1]{\textcolor[rgb]{0.69,0.353,0.396}{#1}}%
\newcommand{\hlkwc}[1]{\textcolor[rgb]{0.333,0.667,0.333}{#1}}%
\newcommand{\hlkwd}[1]{\textcolor[rgb]{0.737,0.353,0.396}{\textbf{#1}}}%
\let\hlipl\hlkwb

\usepackage{framed}
\makeatletter
\newenvironment{kframe}{%
 \def\at@end@of@kframe{}%
 \ifinner\ifhmode%
  \def\at@end@of@kframe{\end{minipage}}%
  \begin{minipage}{\columnwidth}%
 \fi\fi%
 \def\FrameCommand##1{\hskip\@totalleftmargin \hskip-\fboxsep
 \colorbox{shadecolor}{##1}\hskip-\fboxsep
     % There is no \\@totalrightmargin, so:
     \hskip-\linewidth \hskip-\@totalleftmargin \hskip\columnwidth}%
 \MakeFramed {\advance\hsize-\width
   \@totalleftmargin\z@ \linewidth\hsize
   \@setminipage}}%
 {\par\unskip\endMakeFramed%
 \at@end@of@kframe}
\makeatother

\definecolor{shadecolor}{rgb}{.97, .97, .97}
\definecolor{messagecolor}{rgb}{0, 0, 0}
\definecolor{warningcolor}{rgb}{1, 0, 1}
\definecolor{errorcolor}{rgb}{1, 0, 0}
\newenvironment{knitrout}{}{} % an empty environment to be redefined in TeX

\usepackage{alltt}
\usepackage[left=1in,right=1in,top=1in,bottom=1in]{geometry}
\usepackage{amsmath}
\usepackage{hyperref}
\usepackage{float}
%\SweaveOpts[concordance=TRUE]
\IfFileExists{upquote.sty}{\usepackage{upquote}}{}
\begin{document}

\begin{titlepage}
    \centering
    
    { \Huge \textbf{Final Project Report} \par} 
    \vspace{5mm}
    { \huge \textbf{COVID-19's Impact on Motor Vehicle Collisions in NYC} \par} 
    \vspace{3mm}
    { \Large Group 1 (STAT 605 - S1)\par}  
    
    \vspace{5mm}
    \begin{center}
    \line(1,0){425}
    \end{center}
    
    \vspace{5mm}

    \vspace{15mm}
    { \large Authors: \par }
      \vspace{3mm}
    { \Large \textbf{Astha Singh} - \texttt{as677@rice.edu} \par}
    { \Large \textbf{Franco Bosetti} - \texttt{fb57@rice.edu}\par}
    { \Large \textbf{Louis Clarke} - \texttt{lc160@rice.edu}\par}
    { \Large \textbf{Lindsey Russ} - \texttt{ltr1@rice.edu}\par}
    
    
    \vspace{5mm}
    
    \vspace{5mm}

\end{titlepage}

\newpage


\tableofcontents

\newpage

\section{Introduction}

The COVID-19 pandemic created an unprecedented natural experiment in urban mobility. Beginning March 2020, New York City's shelter-in-place orders fundamentally altered how millions of people moved through one of the world's densest transportation networks. This project analyzes how COVID-19 transformed motor vehicle collision patterns across frequency, severity, temporal distribution, and geography.

Using 13 years of police-reported crash data (2012-2025), we examine three distinct periods: Pre-COVID (baseline), COVID (March 2020-December 2021), and Post-COVID (2022-present). Our analysis addresses four questions: (1) Did collision rates return to pre-pandemic levels or establish a "new normal"? (2) Did crashes become more or less deadly during reduced traffic? (3) Were impacts uniform across NYC's five boroughs? (4) How did vehicle characteristics and driver behaviors contribute to collision severity across periods?

\section{Data}

\subsection{Primary Dataset: NYC Motor Vehicle Collisions}
\textbf{Source:} NYC OpenData (2.22M records, 29 variables, 2012-2025)

Police-reported collisions from NYPD's TrafficStat system. Key variables include crash date/time, location (coordinates, borough), casualties (injuries/fatalities by road user type), contributing factors (up to 5 per crash), and vehicle types (up to 5 per crash).

\subsection{Secondary Dataset: Vehicle-Level Details}
\textbf{Source:} NYC OpenData (4.45M records, 25 variables, 2016-2025)

Vehicle-specific data linked via \texttt{COLLISION\_ID}. Each row represents one vehicle in a crash. Variables include vehicle characteristics (make, model, year, type, occupants), driver demographics (sex, license status), pre-crash actions, damage locations, and contributing factors. We aggregated to crash-level metrics: vehicle count per collision, average vehicle age, and average occupants per vehicle.

\subsection{Tertiary Dataset: NYC Speed Limits}
\textbf{Source:} NYC OpenData (141K street segments, Vision Zero initiative)

Posted speed limits by street segment with geographic coordinates. We matched segments to crash locations to analyze whether high-speed roads (>40 mph) correlate with crash frequency and severity changes during COVID. NYC's default speed dropped from 30 to 25 mph in 2014; higher limits indicate highways and major roads.

\subsection{Quaternary Dataset: Twitter Traffic Reports}
\textbf{Source:} Research dataset by Dabiri (2018); 51,178 classified tweets

Tweets categorized as non-traffic (0), traffic incidents (1), or traffic conditions (2). We filtered to 798 NYC-specific tweets, then isolated 66 vehicle crash reports for text mining. While predating COVID, this establishes a baseline for social media reporting patterns and enables contributing factor text analysis from official crash reports.

\subsection{Use of SQL}

In this project, the data preparation process was carried out using a series of SQL queries that transformed raw CSV files into datasets ready for analysis. To keep the workflow organized and easy to follow, each major transformation step was built as a separate SQL view instead of combining everything into one large query. This made the process more modular and transparent, since each view clearly represented a specific stage of data cleaning or filtering. For example, one view focused on identifying tweets mentioning New York City, another detected transit and vehicle references, and a later one isolated road incidents and crashes. Using views in this way made the code easier to read, maintain, and verify, while also allowing intermediate results to be checked throughout the data preparation pipeline.


\newpage












\section{Methods Overview}

\section{Exploratory Data Analysis}

\subsection{Did collision rates return to pre-pandemic levels or establish a "new normal"?}

\subsubsection{Motor Vehicle Collisions over Time}

We begin by examining how the overall number of vehicle collisions in NYC evolved over time. Monthly crash counts show a dramatic decline in early 2020, coinciding with the onset of COVID-19 restrictions and the citywide lockdown. This sharp drop reflects the sudden reduction in traffic volume as commuting, tourism, and social activity came to a halt.

\begin{knitrout}
\definecolor{shadecolor}{rgb}{0.969, 0.969, 0.969}\color{fgcolor}

{\centering \includegraphics[width=\maxwidth]{figure/unnamed-chunk-6-1} 

}


\end{knitrout}

We begin by examining how the overall number of vehicle collisions in NYC evolved over time. Monthly crash counts show a dramatic decline in early 2020, coinciding with the onset of COVID-19 restrictions and the citywide lockdown. This sharp drop reflects the sudden reduction in traffic volume as commuting, tourism, and social activity came to a pause.

When broken down by borough, the same general pattern appears, though the magnitude of change differs. Brooklyn and Queens experienced the largest declines and slowest recoveries, consistent with their heavy commuter traffic. Manhattan saw a steep initial drop followed by a quicker rebound, reflecting its concentration of essential services and commercial traffic. Staten Island and the Bronx showed smaller relative changes, indicating more stable local travel patterns.

Overall, these trends reveal that while traffic volumes rebounded after the lockdown, the city’s collision rates did not fully recover. The data suggest a structural shift in how, when, and why people travel in NYC.

\subsubsection{Yearly Crash Trends by Borough, 2012–2025}

\begin{knitrout}
\definecolor{shadecolor}{rgb}{0.969, 0.969, 0.969}\color{fgcolor}

{\centering \includegraphics[width=\maxwidth]{figure/unnamed-chunk-7-1} 

}


\end{knitrout}

When broken down by borough, the same general pattern appears, though the magnitude of change differs. The figure below shows yearly crash totals by borough from 2012 to 2025. Brooklyn and Queens experienced the largest declines and slowest recoveries, consistent with their heavy commuter traffic. Manhattan saw a steep initial drop followed by a slower rebound, reflecting lasting reductions in commuting and tourism. The Bronx showed a moderate decline followed by gradual recovery, while Staten Island remained relatively stable.

These results indicate that while the overall pattern of decline and partial recovery was shared citywide, the impacts were not uniform. Differences in population density, travel purpose, and road network composition likely explain why some boroughs experienced larger and more persistent changes than others.

\subsection{When do crashes happen? Daily/weekly patterns}

\subsubsection{Crash Frequency by Hour and Day of Week}

To understand when crashes are most likely to occur, we created a heatmap of the crash frequency by hour of the day and day of the week. Each tile represents the total number of crashes in that time slot, with brighter colors indicating more frequent crashes.

\begin{knitrout}
\definecolor{shadecolor}{rgb}{0.969, 0.969, 0.969}\color{fgcolor}

{\centering \includegraphics[width=\maxwidth]{figure/unnamed-chunk-8-1} 

}


\end{knitrout}

The heatmap shows a clear concentration of crashes during daytime and evening hours, especially between 8 AM and 7 PM across all weekdays. Friday stands out with the highest crash intensity, particularly in the late afternoon and early evening, likely reflecting a mix of rush-hour traffic and increased social activity at the end of the work week.

In contrast, the overnight hours (midnight to 6 AM) consistently display the fewest crashes across all days, reflecting lower traffic volumes during those times. Sundays appear somewhat less intense overall compared to weekdays, with a broader distribution of crashes throughout the afternoon.

Overall, the heatmap highlights how collision risk peaks during high-traffic commuting periods and social activity times, with Friday evenings being particularly hazardous in New York City.

\subsubsection{Hourly Crash Distribution by COVID Period}

\begin{knitrout}
\definecolor{shadecolor}{rgb}{0.969, 0.969, 0.969}\color{fgcolor}

{\centering \includegraphics[width=\maxwidth]{figure/unnamed-chunk-9-1} 

}


\end{knitrout}

After examining overall daily and weekly crash patterns, we next examined how these hourly crash patterns evolved across the Pre‑COVID, During‑COVID, and Post‑COVID periods. The proportional distribution of crashes by hour shows that daily patterns remained broadly consistent across the three COVID periods, but the intensity of rush-hour peaks changed. Before the pandemic, collisions followed a clear bimodal pattern with pronounced morning (around 8AM) and evening (5–6PM) peaks. During COVID, these peaks flattened and shifted slightly later in the day, reflecting reduced commuting and more flexible travel times. In the post‑COVID period, the peaks partially recovered but stayed lower than pre‑pandemic levels, suggesting that hybrid work and lasting behavioral changes continue to moderate traditional rush-hour peaks.

\subsection{Did crashes become more severe?}

We next assessed whether crash severity changed during the pandemic by comparing the rates of fatal and injury crashes across periods.

\subsubsection{Nighttime Crashes Are More Severe: Baseline Pattern}

\begin{knitrout}
\definecolor{shadecolor}{rgb}{0.969, 0.969, 0.969}\color{fgcolor}

{\centering \includegraphics[width=\maxwidth]{figure/unnamed-chunk-10-1} 

}


\end{knitrout}

This plot compares the rates of fatal and injury crashes by hour of day to assess whether crash severity changed during the pandemic. The injury rate follows a clear daily pattern, dipping in the early morning hours and rising sharply during daytime traffic peaks. Fatality rates remain much lower overall but show a slight increase overnight, when traffic volumes are low and speeds are higher. As expected, fatalities are far less common than injuries, with the fatality rate remaining close to zero compared to the much higher injury rate.

These patterns suggest that while the total number of crashes declined during the pandemic, the relative risk of severe outcomes did not decrease proportionally. Nighttime crashes continued to be far deadlier, reflecting persistent high‑risk driving behaviors such as speeding and impaired driving.

\subsubsection{Crash Severity by COVID Period}
\begin{knitrout}
\definecolor{shadecolor}{rgb}{0.969, 0.969, 0.969}\color{fgcolor}

{\centering \includegraphics[width=\maxwidth]{figure/unnamed-chunk-11-1} 

}


\end{knitrout}

The plot shows how crash severity shifted across COVID periods. During the pandemic, the fatality rate per 1,000 crashes increased, while the injury rate declined, indicating that although fewer collisions occurred, they were more severe. This pattern suggests that reduced traffic volume led to higher driving speeds and riskier behavior. In the post‑COVID period, injury rates partially rebounded, but fatality rates remained elevated, implying that some high‑risk driving habits persisted even after traffic returned to normal levels. Overall, the plot highlights a lasting change in crash severity linked to pandemic‑era driving behavior.

\subsection{What causes crashes? How did contributing factors change?}

\subsubsection{Causes of Collisions}

To identify the most common causes of collisions, we aggregated all recorded contributing factors across vehicles. The bar chart below shows the five most frequent causes citywide.

\begin{knitrout}
\definecolor{shadecolor}{rgb}{0.969, 0.969, 0.969}\color{fgcolor}

{\centering \includegraphics[width=\maxwidth]{figure/unnamed-chunk-12-1} 

}


\end{knitrout}

Driver inattention or distraction stands out as the leading cause, responsible for more than three times as many crashes as the next factor. Failure to yield, following too closely, and unsafe passing follow far behind, showing that most collisions stem from everyday driving mistakes rather than external conditions.

\subsubsection{Distribution of contributing factors by COVID period}



\newpage

\section{Text Mining (Twitter vs Motor Collision NYC)}

We applied text mining to all five CONTRIBUTING.FACTOR.VEHICLE columns (about 2.2M entries) and to 65+ crash-related tweets. The goal was to compare how Motor Collision NYC data and social media describe traffic incidents.

The comparison in the figure below shows a clear distinction. Motor Vehicle Collision data emphasize behavioral causes such as distraction, failure, and unsafe, while Twitter data focus on location and navigation terms like Bronx, parkway, and street. This suggests that official records capture why crashes occur, whereas social media posts describe where they occur.

\begin{knitrout}
\definecolor{shadecolor}{rgb}{0.969, 0.969, 0.969}\color{fgcolor}

{\centering \includegraphics[width=\maxwidth]{figure/unnamed-chunk-14-1} 

}


\end{knitrout}

We also text‑mined for speed‑related terms across all contributing factor fields to quantify behavioral risk. The table below summarizes the share of crashes mentioning speed.

\begin{knitrout}
\definecolor{shadecolor}{rgb}{0.969, 0.969, 0.969}\color{fgcolor}\begin{table}

\caption{\label{tab:unnamed-chunk-15}Speed-Related Crashes by COVID Period (Text-Mined)}
\centering
\begin{tabular}[t]{l|r|r|r}
\hline
Period & Total Crashes & Speed-Related & Percent\\
\hline
Pre-COVID & 1659997 & 120776 & 7.3\\
\hline
During COVID & 195404 & 18588 & 9.5\\
\hline
Post-COVID & 350217 & 33529 & 9.6\\
\hline
\end{tabular}
\end{table}

\end{knitrout}

Speed‑related mentions increased from 7.3\% pre‑COVID to 9.5\% during COVID and 9.6\% post‑COVID, indicating that risky driving behaviors became more prominent during and after reduced traffic periods. This text‑derived feature reinforces the “severity paradox” observed in structured crash metrics and highlights how unstructured text can reveal behavioral patterns that numeric variables alone may overlook.

\newpage

\section{Statistical Modeling \& Interpretation}

\subsection{Part 1: Baseline Risk Factors for Fatal Crashes}

To investigate how the dynamics of fatal motor vehicle collisions in New York City shifted following the onset of the COVID-19 pandemic, we developed an interaction-based logistic regression model. The outcome of interest remains binary: whether a crash was fatal (at least one person killed) or non-fatal. However, unlike a standard baseline model, this approach allows us to explicitly test whether the relationship between risky behaviors and fatality rates changed significantly between the Pre-COVID and Post-COVID periods.

In simple terms, this model performs a difference-in-differences analysis. It compares crash patterns across thousands of incidents to answer two distinct questions:

\begin{enumerate}
    \item Baseline Risk: Which factors (e.g., motorcycles, speeding) are inherently dangerous? dangerous?
    \item The "COVID Effect": Did these specific factors become significantly more lethal after the pandemic began, even after accounting for general background trends?
\end{enumerate}

To ensure statistical robustness, we aggregated over 50 granular police-reported contributing factors into seven high-level behavioral categories (such as Impairment, Speeding/Aggression, and Distraction). By interacting these categories with the time period and using "Other/Unknown" crashes as a control group, the model isolates distinct behavioral shifts from general systemic changes. We also controlled for borough, time of day, and vehicle type to ensure that observed spikes in fatality risk were due to driver behavior rather than location or vehicle differences.

To interpret the results of the logistic regression, we rely on the Odds Ratio. This metric functions as a risk multiplier, quantifying how strongly a specific factor is associated with a fatal outcome compared to the baseline scenario.

The values are interpreted as follows:

\begin{itemize}
    \item An Odds Ratio of 1.0 indicates no difference in risk. The factor has no impact on the likelihood of a fatality compared to the baseline.
    \item An Odds Ratio greater than 1.0 indicates an increased risk. For example, a value of 3.0 suggests that the presence of this factor makes a fatal outcome three times more likely than if the factor were absent.
    \item An Odds Ratio less than 1.0 indicates a decreased risk. This suggests the factor has a protective effect, making a fatal outcome less likely.
\end{itemize}





While the primary focus of this analysis is the shift in driver behavior, it is crucial to isolate these effects from the inherent physical risks of different crash scenarios. To achieve this, our model controls for Borough, Vehicle Type, and Time of Day as fixed effects. This adjustment accounts for known structural confounders (such as the heightened lethality of motorcycle collisions or the reduced visibility of nighttime driving) which exist independently of the pandemic. By holding these variables constant, we ensure that the significant changes observed in the behavioral categories represent a genuine shift in the lethality of driver conduct, rather than a byproduct of changing traffic composition or location patterns.

\begin{knitrout}
\definecolor{shadecolor}{rgb}{0.969, 0.969, 0.969}\color{fgcolor}

{\centering \includegraphics[width=\maxwidth]{figure/unnamed-chunk-18-1} 

}


\end{knitrout}

This plot illustrates the fixed effects of the control variables (Vehicle Type, Time of Day, and Borough), holding all behavioral factors constant. These results validate the model’s reliability, as they align with established traffic safety principles. The most significant predictor of fatality is vehicle type: incidents involving motorcycles show a drastically elevated risk (Odds Ratio > 3.0) compared to the baseline, confirming that motorcyclists remain the most vulnerable road users regardless of other conditions. Time of day is also a significant factor, with nighttime crashes showing nearly double the odds of a fatal outcome compared to daytime incidents, likely due to reduced visibility and higher travel speeds. Geographically, Manhattan is associated with the lowest fatality risk, likely a result of congestion and lower average speeds. On the other hand, Staten Island, with its more suburban, highway-centric infrastructure, presents a higher baseline risk.



\begin{knitrout}
\definecolor{shadecolor}{rgb}{0.969, 0.969, 0.969}\color{fgcolor}

{\centering \includegraphics[width=\maxwidth]{figure/unnamed-chunk-20-1} 

}


\end{knitrout}

This figure presents a comprehensive view of how driver behaviors shifted during the pandemic. The left panel illustrates the absolute risk levels, identifying Speeding/Aggression and Traffic Violations as the most lethal factors in both periods (Odds Ratios > 2.0).

However, the right panel reveals crucial nuances in how these risks changed. Notably, Distraction and Traffic Violations showed a statistically significant decrease in lethality relative to the baseline trends (Interaction OR < 1.0, $p < 0.05$). In contrast, while Impairment appeared to become more dangerous in the raw data (Left Panel, Red dot > Blue dot), the interaction model (Right Panel, Gray bar) confirms that this shift was not statistically distinguishable from chance ($p > 0.05$). This suggests that while distracted driving became less predictive of fatalities, the extreme danger of impaired driving remained a stable constant throughout the pandemic.

\subsubsection{Interpretation}

In summary, this analysis utilized an interaction-based logistic regression to disentangle the specific effects of the COVID-19 pandemic from established structural risk factors in NYC collisions. By controlling for fixed effects such as vehicle type and borough, we confirmed that physical determinants of fatality—most notably motorcycle involvement and nighttime driving—remained consistent predictors of mortality regardless of the pandemic. 

However, the interaction model revealed that the "COVID Effect" was not a uniform amplifier of risk. While Speeding/Aggression and Impairment remained persistently high-risk behaviors with no statistically significant deviation from their baseline lethality, Distraction and Traffic Violations exhibited a significant relative decline in their association with fatal outcomes. These findings underscore that the post-pandemic traffic landscape was not defined by a blanket increase in danger across all categories, but rather by a specific restructuring of how certain driver behaviors translate into fatal outcomes.


\subsection{Part 2: COVID-19's Impact on Crash Severity}

Building on our baseline risk analysis, we developed statistical models to quantify how COVID-19 specifically altered crash severity patterns in NYC. This analysis tests whether the pandemic period itself was independently associated with changes in fatality risk, controlling for established risk factors.

\begin{knitrout}
\definecolor{shadecolor}{rgb}{0.969, 0.969, 0.969}\color{fgcolor}

{\centering \includegraphics[width=\maxwidth]{figure/unnamed-chunk-21-1} 

}


\end{knitrout}

The statistical models reveal that COVID-19 had a significant and substantial effect on crash severity. In the unadjusted model, crashes during COVID were 1.91 times more likely to be fatal (95\% CI: [1.67, 2.19]), and post-COVID crashes were 2.12 times more likely (95\% CI: [1.92, 2.34]) compared to the pre-pandemic baseline.

Even after controlling for differences in borough, vehicle type, and time of day, the effect persisted: crashes during COVID remained 1.67 times more likely to be fatal (95\% CI: [1.45, 1.91]), and post-COVID crashes were 1.79 times more likely (95\% CI: [1.61, 1.98]). The persistence of this strong effect in the adjusted model suggests the pandemic introduced fundamental changes in crash dynamics beyond simple compositional shifts in traffic patterns.

\begin{knitrout}
\definecolor{shadecolor}{rgb}{0.969, 0.969, 0.969}\color{fgcolor}

{\centering \includegraphics[width=\maxwidth]{figure/unnamed-chunk-22-1} 

}


\end{knitrout}

\subsubsection{Interpretation}

The interaction analysis reveals that while overall severity increased across all vehicle types, the relative risk hierarchy remained largely consistent. No individual vehicle type showed statistically significant differential effects during COVID at the 0.05 level, though motorcycle crashes showed a notable trend toward increased risk during COVID (p = 0.097) and truck/bus crashes trended toward increased risk post-COVID (p = 0.100). This suggests the pandemic amplified crash severity uniformly rather than disproportionately affecting specific vehicle categories.

\newpage

\section{Killer Plot}

\newpage

\section{Conclusion}

\subsection{Key Takeaways}

\subsection{Potential Future Work}

\newpage

\section{Appendices}

\subsection{Appendix A: Piecewise Linear Trends in Total Crashes per Year by Borough}
\begin{knitrout}
\definecolor{shadecolor}{rgb}{0.969, 0.969, 0.969}\color{fgcolor}

{\centering \includegraphics[width=\maxwidth]{figure/unnamed-chunk-23-1} 

}


\end{knitrout}

\subsection{Appendix B: Total Number of Accidents per Time of Day}
\begin{knitrout}
\definecolor{shadecolor}{rgb}{0.969, 0.969, 0.969}\color{fgcolor}

{\centering \includegraphics[width=\maxwidth]{figure/unnamed-chunk-24-1} 

}


\end{knitrout}

\subsection{Appendix C: All Accidents}
\begin{knitrout}
\definecolor{shadecolor}{rgb}{0.969, 0.969, 0.969}\color{fgcolor}
\includegraphics[width=\maxwidth]{figure/unnamed-chunk-25-1} 
\end{knitrout}

\subsection{Appendix D: Highway Accidents}
\begin{knitrout}
\definecolor{shadecolor}{rgb}{0.969, 0.969, 0.969}\color{fgcolor}

{\centering \includegraphics[width=\maxwidth]{figure/unnamed-chunk-26-1} 

}


\end{knitrout}

\subsection{Appendix E: R Code}
\begin{knitrout}
\definecolor{shadecolor}{rgb}{0.969, 0.969, 0.969}\color{fgcolor}\begin{kframe}


{\ttfamily\noindent\bfseries\color{errorcolor}{\#\# Error in parse(text = input): <text>:1:4: unexpected symbol\\\#\# 1: \{r ref.label\\\#\# \ \ \ \ \ \ \ \textasciicircum{}}}\end{kframe}
\end{knitrout}

\end{document}
